\section{Dataset Quality}

In this section, we briefly review a few properties that must be satisfied by samples contained within a dataset in order to ensure that systems trained on it do not overestimate their performance on unseen data as a result of poor generalization.
A more complete discussion of this topic can be found in Croft et al.\cite{croft2023data} which deals with datasets for software vulnerability detection.
The data quality attributes discussed by Croft et al. are standardized in the ISO/IEC 25012 data quality framework\cite{iso/iec25012}.
They employ these attributes in an analysis of several datasets including Big-Vul.
Let's consider the data quality attributes relevant to our work.

\subsection{Uniqueness}
A datum is \emph{unique} if it is represented exactly once within the dataset.
Supervised learning systems typically employ datasets that are partitioned into one or more of training, validation, and test sets.
Duplicates within any of these subsets can lead to biased model performance.
Duplicates across these subsets lead to data leakage which can have the negative result of overestimating model performance.

\subsection{Consistency}
In the context of software vulnerability repair datasets, a datum is \emph{consistent} if it does not appear in the dataset with conflicting labels.
There are two possible reasons for a single code sample to be labeled in different ways in a data set:

\begin{enumerate}
    \item The labels employed do not partition the space of possible data, that is, they are not exclusive of each other.

    \item At most one of the labels for such a sample is the appropriate label and all other labels are incorrect.
\end{enumerate}

Inconsistent labels arising from either of these causes are inappropriate for machine learning systems intended to identify or repair security vulnerabilities.
If the system is to identify the unique class associated with a program vulnerability, it must be trained using samples with only a single label.
On the other hand, if the system is to repair such vulnerabilities, then inconsistent labels will make it more difficult to learn class-specific repairs. 

\subsection{Accuracy}

A datum in a dataset is \emph{accurate} if it embodies \emph{semantic label correctness}.
In the setting of a vulnerability dataset aimed at program repair, this means samples must correctly identify code as being vulnerable or correct.
If the dataset identifies the particular vulnerability exhibited by code, that vulnerability must be present in the sample.
If the dataset marks the vulnerable code, it must mark it correctly.
If a repaired sample is provided, the vulnerability must be correctly repaired.
And if the dataset marks the repaired code, it must mark it correctly.

Numerous samples in the VulRepair dataset are labeled with the wrong CWE tag.
This occurs for a variety of reasons.
In some cases, a change appears in a commit that addresses a CVE but the specific change is unrelated to that CVE.
In other cases the CWE labels from Big-Vul and CVEfixes do not agree as a result of a change to the NIST database during the time between the creation of the two data collections.

In some cases, correct code is labeled as being incorrect.
Other samples contain incomplete target patches.
In addition to such labeling problems, there are meta-problems associated with the sample labeling approach adopted by these datasets.
One such case is samples in which code must be moved in order to address a vulnerability.
Applying \texttt{<Startbug>} and \texttt{<EndBug>} labels to such incorrectly placed code is conflated with the use of these tags to identify code that must be changed in-place.
Perhaps we need an alternate mechanism to identify repairs that involve code motion rather than modification.




\subsection{Completeness}
A vulnerability repair dataset sample is considered to be \emph{complete} if the sample itself contains all information necessary to determine that the identified vulnerability exists and is correctly repaired.

\subsection{Discussion}
\label{Dataset:Discussion}

When Fu et al. published their work on VulRepair, their dataset was assigned two different quality certifications by ACM, namely, \emph{Artifacts Available} and \emph{Artifacts Evaluated - Reusable}\cite{acm-badges}.
The second of these badges is to be assigned only to artifacts that have successfully completed ``an independent audit.''
The nature of the audit is not reported, but clearly that audit did not effectively assess the data quality attributes discussed above.

Construction of datasets containing the numbers of samples necessary to effectively fine-tune code language models to deal with security vulnerability repair is a difficult task.
Attempts to construct such data sets in an automated way have been fraught with data quality problems.
We discuss our recommendations concerning construction and maintenance of such datasets in Section \ref{Conclusion}.